\begin{table}[H]
\centering
\begin{threeparttable}
\caption{Construction of the Three Headline Summary Indices}
\label{tab:index_construction}
\scriptsize
\setlength{\tabcolsep}{4pt}
\begin{tabularx}{\textwidth}{p{3.2cm} p{1.2cm} X X}
\toprule
Index & Wave(s) & Components & Construction \\
\midrule
Legitimacy index & Endline, follow-up & Regime support index; censorship support index; trust in state media; trust in the Communist Party; trust in society & Unit-weighted average of standardized components. Each component is standardized using the randomized sample\textquotesingle s baseline distribution so that endline and follow-up changes are on a common pre-treatment scale. \\\\
\addlinespace
Foreign-media receptivity-and-demand index & Endline & Trust in foreign media; mean weekly credibility; mean weekly desire for similar content; willingness to pay bid; purchase indicator & Unit-weighted average of standardized components. Each component is standardized using the endline control-group distribution so the index is centered on the untreated post-course benchmark for comparable respondents. \\\\
\addlinespace
Identity-backlash index & Endline & Nationalism index; nationalism-bias item; reverse of trust in foreign media & Unit-weighted average of standardized components. Each component is standardized using the randomized sample\textquotesingle s baseline distribution, with trust in foreign media reverse-coded so larger values indicate stronger backlash. \\\\
\bottomrule
\end{tabularx}
\begin{tablenotes}[flushleft]
\footnotesize
\item Note: These headline indices are not principal-component-analysis (PCA) scores or item-response-theory (IRT) scores. They are pre-specified, unit-weighted standardized averages designed to summarize substantively aligned but not identical outcomes. The underlying attitudinal batteries have conventional scale diagnostics reported in Table~\ref{tab:scale_reliability}; the foreign-media receptivity-and-demand index is intentionally a design-based composite spanning beliefs, weekly evaluations, and consequential choices.
\end{tablenotes}
\end{threeparttable}
\end{table}
